% Standalone source-record documentation; no research attempts.
\documentclass[11pt]{article}
\usepackage[margin=25mm]{geometry}
\usepackage{fontspec,unicode-math}
\setmainfont{Latin Modern Roman}
\setmathfont{Latin Modern Math}
\usepackage{amsmath,amssymb,mathtools,xurl,hyperref}
\hypersetup{colorlinks=true,urlcolor=blue,linkcolor=blue}
\setcounter{secnumdepth}{0}
\setlength{\parindent}{0pt}
\setlength{\parskip}{5pt}
\setlength{\emergencystretch}{3em}
\newcommand{\sref}[2]{\hyperlink{source:#1:#2}{[S#2]}}
\newcommand{\cataloguescope}[1]{\paragraph{Source record.}#1}
\newcommand{\floor}[1]{\left\lfloor#1\right\rfloor}
\begin{document}
% UnsolvedMath ID 20002205; source catalogue code AIM-INFRASTRUCTURE-0105
\setcounter{section}{20002204}
\section{Identification and repair of an AB/BA proof-learning study}\label{20002205}
{\small\sffamily UnsolvedMath ID 20002205 · Computer Science}
\subsection{Definitions and mathematical statement}

\paragraph{Shared notation.}Standard number systems are $\mathbb N,\mathbb Z,\mathbb Q,\mathbb R,\mathbb C$. Specialized notation and index ranges follow the cited source and the definitions below.


Inquiry-based learning (IBL) engages students in mathematical investigation and justification. A real-analysis course studies topics such as limits, continuity, and proofs from definitions. A classroom culture consists of expectations and practices for participation, explanation, feedback, and responsibility. Proof production asks a student to construct an argument; proof comprehension concerns understanding an argument that is presented. An order comparison distinguishes attempting a proof before viewing it from viewing it before attempting it. Crossed conditions can also distinguish a learner's usual classroom setting from a particular presentation condition.
\paragraph{Record scope and qualifications.}
The catalogue title is retained as metadata. Named certificates, contracts, bounds or frameworks in that title are not asserted as results of the original source extract.
The long extract includes additional developmental and course-study proposals. Measures, populations, instructional implementations, and assumptions for comparison must be specified; described studies are proposals, not verified results.

Catalogue screening limitation: Empirical education study and classroom-culture questions, rather than a mathematical statement.

The source snapshot labels this record partially solved; this is catalogue metadata rather than independent verification of a mathematical result or the current completion of a project.
The source text below is rendered literally where necessary; damaged formulas and composite extracts are not silently repaired into a new mathematical claim.
The supplied catalogue marks this formulation as an unverified reconstruction; missing assumptions are not certified by the source link.
\paragraph{Statement recorded in the supplied source.}
(1) How does students having attempted a proof production influence their learning from a proof presentation? (2) How do students ongoing classroom experiences (in lecture or IBL classrooms) influ-ence their learning from proof production and observation of a proof presentation? 4

Abstract: IBL instructors anticipate that students learn from their peers proof presenta-tions in part because they have already attempted the task themselves. To investigate this, we propose that two groups of students attempt to (a) produce a proof of a claim and (b) view a presentation of a proof of that claim, in alternate orders. The condition of proof production before presentation mimics the IBL listening environment. After both learning experiences, Mejia-Ramos et al.s proof comprehension instrument will be administered to assess various dimensions of student learning about the proof. We further propose that stu-dents currently enrolled in both IBL and lecture-style courses participate in the study. This will make the study sensitive to the possibility that students ongoing learning practice may determine their optimal learning conditions (rather than the learning conditions themselves).

• Alternative possible conditions: proof presentation by an experienced instructor or proof presentation by a student (with or without mistakes).

• The study design may need to attend to how students time listening to their peers proof presentations is structured by the instructor (e.g. by assigning roles, by pro-viding rubrics, by inviting peer-peer feedback).

• We anticipate that the choice of mathematical topic, proof task, and that tasks relative difficulty will be crucial to the outcome and viability of the study.

Research Question: How does student learning from a proof presentation differ between presentations made by expert mathematics faculty and presentations made by other students?

Abstract: Students are the primary authors of mathematical proofs in IBL courses, meaning that students will more often see imperfect proofs rather than valid proofs that exhibit standard mathematical form. Observing peer proofs may invite more active listenership since students are expected to question and give feedback on peer proofs. Observing expert proofs may improve learning because they serve as models of appropriate proof writing. We propose a study that compares students proof comprehension, as measured by Mejia-Ramos et al.s instrument, after viewing the two types of proof presentations. We anticipate that conducting the study with students from both IBL and lecture-style courses will benefit the study. It will be important to distinguish the comparative benefits of presentation conditions from the possibility that students simply learn how to listen effectively within their native learning environment.

• This study was conceptualized with reference to similar studies in physics education and elsewhere where listening to peers produced greater gains in comprehension over a short time interval relative to some measures of learning.

• We anticipate that the 2x2 design coupled with the multi-dimensional measure of proof comprehension will provide a rich set of possible outcomes leading to different inferences about the nature of student learning from listening.

Persistence \& identity. Exploratory study of student development.

Research Questions:

• What are the student developmental categories in a Moore Method course?

• How can we refine and explain these categories?

• How can we explain and describe student advancement between developmental cat-egories? 5

Abstract: Many Moore Method 1 instructors report seeing dramatic transformations ex-perienced by students when they make significant steps forward in their levels of mathe-matical achievement. Through qualitative observation of several Moore Method classrooms, we anticipate identifying and refining definitions of student developmental categories. We anticipate this study will develop baselines for subsequent studies to document transforma-tions to higher developmental categories in Moore Method classes, record them in detail, and investigate the conditions that foster such transformations. We propose a preliminary list of developmental categories as follows:

Beginners:: Students who have not been able do any presentations successfully.

Novices:: Students who successfully present a proof that is a follow your nose type proof (they know what a definition is and definitions are and the logic and format of a proof are and can put them together).

Apprentices:: Students who successfully present a proof that requires some significant insight or ingenuity to answer.

Masters:: Students who initiate interest in mathematical problems that they want to do themselves. Note that these levels somewhat parallel Lee May's levels of performance, but differ in a practical way. Assessing a student for May's levels require a carefully designed assessment procedure in addition to the classwork itself. In contrast, instructors can classify students according to the levels proposed here based only on what they see from the student's class performance.

Data Sources: Ted Mahavier has video we might use for exploratory studies of his and his father's courses. Anneliese Spaeth and Padraig McLoughlin have volunteered to use their Analysis courses and possibly record them in Spring 2016. Ted Mahavier is on sabbatical during Spring 2016, and has volunteered some time. Ted Mahavier is teaching Real Analysis in Fall 2016, and has volunteered his course for participation in this study. We invite others to participate by contributing observations of these transformations in their own Moore Method classes.

Case study of the impact of IBL on student development.

Research Questions: How does taking a Moore Method course affect student:

• Confidence: Willingness to engage and belief in eventual success

• Independence: Prove theorems and solve problems (and verify on their own)

• Identity: Participation in mathematical culture

• Willingness to try and fail (and see its worth)

• Resilience: Willingness to try after failure

• Perception of self worth/worth of their work

• Locus of control

• Perception of the nature of mathematics

> 1By Moore Method, we mean an Inquiry Based Learning (IBL) course centered around individual student presentations. 6

Abstract: The Colorado Study (2014) 2 established that students from IBL courses attain greater success in subsequent math courses than students from more traditional, lecture-based courses. In an effort to explain in more depth why this might be the case, we propose a study tracking any of several student attributes identified by mathematicians as important for student success in mathematics. Above we have begun a preliminary description of these attributes. We propose a series of case studies of students making transitions from one category (Beginner, Novice, Apprentice, Master) to another, including cases of failure to transition.

Benefits of MMM over Lecture for Strong Students.

Research Questions: \sref{20002205}{2}
\paragraph{Source status.}
The supplied catalogue labels this source record partially solved; that metadata was checked against the supplied snapshot on 2026-10-11. The label is preserved for provenance and is not an independent verification of present open status or mathematical correctness.
\subsection{Short English statement}
Investigate learning from proof production and presentations, including the order of activities and the presenter's role.
\subsection{Sources}
\begin{itemize}
\item[\hypertarget{source:20002205:1}{S1}] Ulam AI and the UnsolvedMath Contributors, UnsolvedMath: A Curated Collection of Open Mathematics Problems, record 20002205 (AIM-INFRASTRUCTURE-0105). Catalogue curation and metadata: CC BY 4.0. The source extract is retained; definitions, English summaries and qualifications are editorial additions.
\url{https://huggingface.co/datasets/ulamai/UnsolvedMath}.
\item[\hypertarget{source:20002205:2}{S2}] AIM Workshop Problem Lists, source linked by catalogue record AIM-INFRASTRUCTURE-0105.
\url{http://aimath.org/pastworkshops/iblanalysisrep.pdf}.
\end{itemize}
\paragraph{Preparation and review.}
The original source and UnsolvedMath contributors retain attribution for the source extract. Codex (GPT-6; exact serving revision not exposed) assisted with definitions, English summaries, scope notes and formatting. No new human mathematical authorship or independently verified solution is asserted. Permanent source ID 20002205 and code AIM-INFRASTRUCTURE-0105 checked against the complete supplied registry and current repository. Distinct upstream records, including repeated questions and broad or damaged extracts, are retained for complete coverage of the source CS category. A repeated record is not a new independent conjecture. Complete mathematical-equivalence and current-literature review has not been performed. Preparation and final review follow contribution prompts C and D, with source gaps disclosed. Structural and rendering checks do not establish a new theorem or open status.
\label{end:20002205}
\end{document}
